\scanpage{30}
\mspage{8}
\margin{29}

(ii) L'homomorphisme canonique
\[
\deleted{H^i(X,F)}\ R^i\widehat f(\widehat F)\qquad \lim R^if(F_m)
\]
est un isomorphisme pour $i\quad n$, un monomorphisme pour $i=n+1$ \deleted{\unsure{;}}

Démonstration. On se ramène aussitôt au cas où $S$ est affine, et à prouver alors le

\uline{Corollaire} 2.2. Sous les conditions de 2.1. supposons de plus $S$ affine.
Alors :

(i) L'homomorphisme canonique
\[
H^i(X,F)\qquad H^i(\widehat X,\widehat F)
\]
est un isomorphisme pour $i\quad n$, un monomorphisme pour $i=n$.

(ii) L'homomorphisme canonique
\[
H^i(\widehat X,\widehat F)\qquad \lim H^i(X_m,F_m)
\]
est un isomorphisme pour $i\quad n$, un monomorphisme pour $i=n+1$.

\deleted{Pour \unsure{Démontrer|démontrer}} Quitte à remplacer $\underline O_X(1)$ par une puissance tensorielle, et $t$ par une puissance de $t$, on peut supposer $\underline O_X(1)$ très ample relativement à $S$. D'autre part, $t$ donc $t^m$ étant $F$-régulier, la multiplication par $t^m$, considéré comme homomorphisme de $F(-m)$ dans $F$, est injectif, donc on a pour tout $m\quad 0$ une suite exacte :
\diagram{\begin{tikzcd}[column sep=large]
(14) & 0 & F(-m)\arrow[r,phantom,"t^m"] & F & F_m & 0,
\end{tikzcd}}
d'où une suite exacte de cohomologie
\diagram{\begin{tikzcd}[column sep=large]
H^i(X,F(-m)) & H^i(X,F) & H^i(X,F_m) & H^{i+1}(X,F(-m)).
\end{tikzcd}}
Or en vertu de 1.5. on a $H^i(X,F(-m))=0$ pour $i\quad n$, et $m$ grand, ce qui prouve le

\uline{Lemme} 2.3. \deleted{\unsure{On}} Pour $m$ grand, l'homomorphisme canonique
\[
H^i(X,F)\qquad H^i(X_m,F_m)
\]
est bijectif si $i\quad n$, injectif si $i=n$.

Cela montre que pour $i\quad n$, le système projectif $(H^i(X_m,F_m))_m$

